% Self-contained experimental section. Requires amsmath, amssymb, graphicx,
% booktabs, hyperref. Define \experimentroot if this file is included elsewhere.
\providecommand{\experimentroot}{.}
\section{Completed audits as future fitting data}
\label{sec:audit-release-diagnostic}

The following new-seed diagnostic asks whether permanently excluding completed
audits from later fitting creates avoidable limitations in a trained adaptive
learner. It does not claim a new testing theorem, and it does not implement
transfer of historical certification evidence to a changed prediction pair.
The broad motivation is useful structural adaptation under explicit data and
computation costs; success of the residual proposal is not assumed.

\subsection{Implemented evidence mechanism}
Let $H_{t-1}$ contain every earlier fitting/audit observation, decision,
parameter and optimizer state. Current proposal randomness may be included
before the next audit. Suppose all current predictors and comparison margins
are then fixed and the current block $A_t$ is independent and iid. For old
minus new bounded-loss differences $D_i\in[-1,1]$, define
\[
 E_n=\frac15\sum_{\lambda\in\Lambda_\tau}
       \prod_{i=1}^n\{1+\lambda(D_i-\tau)\},\qquad
 \Lambda_\tau=\{.05,.1,.2,.5,(1+\tau)^{-1}\},
\]
with the implemented margins $\tau\in\{0,.0005\}$.
Accept the inequality when $E_n\ge 1/\alpha_{t,j}$. This is an established
bounded-mean betting construction, not a new test; the direct foundation is
Waudby-Smith and Ramdas, Section~4 and Appendix~A.3 of the inspected
arXiv version~7 \cite{waudby2024}.

Each product has nonnegative factors and conditional expectation at most one
under its null. Its fixed average therefore gives the conditional error bound
needed by the preceding proposition. Every changed prediction pair starts a
new product. The current experiment uses fixed blocks, with no uncorrected
multiple inspections within a block.

\subsection{Frozen design and actual execution}
The protocol and runner hashes were frozen at 03:45:05 UTC on
15 September 2026, before seeds 750000--750011 were run. Four tasks and
12 independent seeds per task give 48 paired configurations, seven methods
per configuration, 336 complete trajectories and 8,064 events. All completed
the preset 24-event horizon. The main run took 330.7 seconds with two local
CPU workers. An unchanged-runner recovery of the earlier interaction seed
63000 matched its full two-method history and controls before this study.

Write $\sigma(u)=1/(1+e^{-u})$. All tasks have $X\sim N(0,I_6)$ and
$Y\mid X\sim\operatorname{Bernoulli}(.05+.9\sigma(g(X)))$, where
\begin{align*}
g_{\rm null}(x)&=0,\\
g_{\rm additive}(x)&=2.5\tanh(2x_1+.6)-2\tanh(2x_2-.5)
 +1.7\tanh(2x_3+.3)-1.5\tanh(2x_4),\\
g_{\rm interaction}(x)&=4\tanh(2x_1)\tanh(2x_2)
 +2\tanh(2x_3)\tanh(2x_4),\\
g_{\rm radial}(x)&=3(x_1^2+x_2^2-1.4).
\end{align*}
The radial boundary is an additional task with four nuisance coordinates,
not an instance of a supplied-module theorem. The fitted network is
$p_\theta(x)=\sigma(c+\sum_{j=1}^k a_j\tanh(w_j^\top x+b_j))$ with Brier
loss; all its parameters train. There are 2,048 initial fitting labels and
4,096 fresh audit labels at each event, hence 100,352 unique observed labels
per method. Conditional label noise is integrated on 32,768 independent
evaluator inputs per seed; the risk estimate remains Monte Carlo in $X$.
The evaluator is called only after an entire method's action sequence is
fixed and is absent from the learner's arguments.

Four adaptive arms cross withheld/released audit data with residual/random
proposals. The residual score is
$(\overline{rz})^2/(\overline{z^2}+10^{-3})$, where
$r=y-p$ and $z=p(1-p)\tanh(w^\top x+b)$. Four incoming-vector candidates
receive 40 minibatch ascent updates, and one is selected on at most 4,096
historical fitting points. Random additions use the same initial candidate
distribution and random selection. Both add a unit with outgoing weight
zero, preserving the current predictor before branch fitting. Each event
fits both continuation and growth from the incumbent. Growth is accepted
only after its gains over both incumbent and continuation exceed .0005;
otherwise a certified positive continuation is accepted, or the incumbent
is held exactly.

Three controls receive the same label arrivals and release times: persistent
online learners at fixed widths two and 26, and a fixed width-26 learner
using a fresh gate on continuation. Width 26 equals the adaptive maximum
$2+24$ and is specified before outcomes. The ungated controls make no
checkpoint-monotonicity claim. The adaptive arms allocate $.05/(3\cdot24)$
to each comparison and the one-comparison fixed gate $.05/24$, both spending
the same trajectory error budget. That budget is not a simultaneous
guarantee across experimental methods and seeds.

All arms receive six million parameter--example updates for warmup and
eight million per event, including candidate search and rejected branches,
with small integer remainders recorded. They use persistent Adam
$(.9,.999)$, learning rate .01, denominator constant $10^{-8}$ and minibatches
of at most 256. Proposal ascent uses learning rate .04. Branches clone the
incumbent's optimizer state; only the selected branch commits it. New-unit
moments start at zero and retain the existing step counter. Random proposals
allocate saved search work to branch fitting. Fixed learners devote the
entire event budget to continuation. The 198-million-update horizon and
32-million-update ungated tail match an operation proxy, not FLOPs, optimizer
steps, memory or wall time. These shared persistent-minibatch conventions
differ from the earlier full-batch/reset runner; the new factorial comparison
isolates release within the new protocol.

\subsection{Observed release effects and negative controls}
Table~\ref{tab:release-effects} reports primary paired effects. Confidence
intervals are descriptive 95\% Student-$t$ intervals across the 12 independent
seeds; there is no multiplicity adjustment. Event time points and proposals
are not treated as independent replications.

\begin{table}[tb]
\centering\scriptsize
\begin{tabular}{lrrr}
\toprule
Task & Release: residual & Release: random & Residual minus random\\
\midrule
Null & $-1.994\;[-2.813,-1.176]$ & $-2.016\;[-2.846,-1.186]$ & $0.022\;[-0.054,0.097]$\\
Additive & $-4.576\;[-6.648,-2.505]$ & $-3.943\;[-5.666,-2.221]$ & $0.063\;[-0.237,0.363]$\\
Interaction & $-3.721\;[-4.557,-2.885]$ & $-4.657\;[-6.277,-3.038]$ & $-0.110\;[-0.744,0.524]$\\
Radial & $-1.142\;[-2.131,-0.152]$ & $-1.199\;[-2.081,-0.318]$ & $0.554\;[-0.171,1.280]$\\
\bottomrule
\end{tabular}
\caption{Final Brier-risk differences, in units of $10^{-3}$, with descriptive
95\% paired intervals. ``Release'' means released minus withheld. The last
column compares proposals under release. Negative values favor the first
method. All four release-by-proposal interaction intervals also contain zero.}
\label{tab:release-effects}
\end{table}

Both proposal mechanisms benefit from release in the four measured means.
Their released comparison does not resolve a residual-proposal advantage;
intervals containing zero do not establish equivalence. Mean released
residual final risks are .250716, .131908, .144304 and .121427 on null,
additive, interaction and radial, with respective widths two, four, six and
three. Permanently withheld residual risks are .252710, .136484, .148025 and
.122568. Improvement therefore need not be accompanied by additional
growth: mean interaction width remains six under both release policies.

The fixed width-26 online learner has final mean risks .250857, .131491,
.137553 and .117145. Released residual minus this control is
$+.006752\;[.006203,.007300]$ on interaction and
$+.004282\;[.003656,.004907]$ on radial. Even the gated fixed architecture
outperforms released residual in those comparisons:
$+.003801\;[.003225,.004378]$ and
$+.002762\;[.002014,.003509]$, respectively. Thus removal of the gate is not
the sole explanation for these fixed-capacity results. Fixed width two is
strong on null (mean .250107) but underfits interaction (.209650) and radial
(.192170). The larger control is needed for credible attribution.

After the preset ungated tail, released residual risks become .250161,
.131254, .142542 and .118292. The release-minus-withheld effects remain
negative, but the interaction/radial gaps to fixed26 online remain
$+.004743\;[.004021,.005464]$ and
$+.001445\;[.001061,.001829]$. On additive, this post-training comparison
instead favors the smaller released residual network by
$-.000329\;[-.000491,-.000167]$. These tails are diagnostics, not protected
deployment updates. All adverse and favorable task outcomes are retained.

\begin{figure}[tb]
\centering\includegraphics[width=\linewidth]{\experimentroot/figures/risk_trajectories.pdf}
\caption{Actual 24-event histories; curves are seed means and bands are
$\pm1$ standard error. Released labels improve the adaptive arms, while
fixed-capacity controls expose both underfitting at width two and stronger
risk performance at width 26. All arms share data arrivals and optimization
proxy budgets.}
\label{fig:release-trajectories}
\end{figure}

\subsection{Resource ledger and verification}
Only 96,256 labels have been released before fitting the final online
candidates; the last audit is released after that decision. Withheld arms
fit on exactly 2,048 distinct labels. Released adaptive arms fit on roughly
96,061--96,256, while fixed26 fits about 90,625--90,644 within its matched
parameter--example budget. Repeated fit accesses still differ: released
residual averages 9.84m, 6.36m, 4.71m and 7.31m across tasks, versus .947m
for fixed26 and 11.647m for fixed2. Pool availability is not silently
equated with actual fitting access.

Residual candidate optimization costs 6.881m parameter--example updates over
the horizon. Rejected branch fits alone average 168.727m--182.226m for
released residual. Audit prediction work adds 5.800m--14.090m
parameter--example evaluations, versus 41.091m for fixed26 gated. Adaptive
gates perform 1,474,560 betting-factor evaluations per trajectory; fixed26
gated performs 491,520. Prediction and selection work is logged separately.
The adaptive networks have substantially lower inference capacity than
fixed26, so lower fixed-model risk is a resource tradeoff rather than a
universal dominance statement.

Independent reconstruction from the raw data and saved branch parameters
checks all 8,064 events, 14,976 betting comparisons and evaluator risks;
maximum log-e discrepancy is $9.95\times10^{-13}$. All 48 run hashes,
336 access ledgers and four frozen files pass integrity checks. A separate
full seven-arm interaction replay matches all 168 events at relative
tolerance $10^{-10}$ and absolute tolerance $10^{-12}$, excluding elapsed
times. Checks also verify gradients, function-preserving additions, frozen
candidate hashes, immutable commitments, post-decision release and
current/future-label perturbations. No gated checkpoint increases the
evaluator's estimated risk in this sample; that observation is not the
conditional proof.

\begin{figure}[tb]
\centering\includegraphics[width=\linewidth]{\experimentroot/figures/paired_diagnostics.pdf}
\caption{Prespecified paired final-risk contrasts. Bars are descriptive
95\% intervals across independent seeds, without multiplicity correction.
The scale is $10^{-3}$ Brier-risk units.}
\label{fig:release-paired}
\end{figure}

\begin{figure}[tb]
\centering\includegraphics[width=\linewidth]{\experimentroot/figures/risk_capacity_tail.pdf}
\caption{Risk and retained capacity. Filled markers show final committed
mean risks with 95\% intervals; crosses show means after the additional
ungated tail. The fixed26 risk advantage on harder tasks comes with larger
inference capacity.}
\label{fig:release-capacity}
\end{figure}

\paragraph{Scientific status.}
The study diagnoses avoidable permanent withholding under a controlled
learner and supplies same-arrival online controls. It does not demonstrate a
new information score, retained certification evidence, optimized resource
allocation or a practical deep-learning advantage. Synthetic scope,
finite horizons, constant optimizer rates and proxy-compute matching remain
limitations. Complete event and run summaries, a representative JSON/NPZ trajectory,
commands, the frozen protocol, and plotting/audit code are in the repository.
Other raw JSON/NPZ outcomes are listed by hash in the local-archive manifest. It is an
experimental component of the evidence/learning investigation.

\begin{table}[p]
\centering\small
% Generated from actual confirmation run_summary.csv.
\begin{tabular}{llrrr}
\toprule
Task & Method & Final risk & Tail risk & Width \\
\midrule
Null & Withheld residual & 0.25271 & 0.25410 & 2.00 \\
Null & Released residual & 0.25072 & 0.25016 & 2.00 \\
Null & Withheld random & 0.25271 & 0.25410 & 2.00 \\
Null & Released random & 0.25069 & 0.25015 & 2.00 \\
Null & Fixed 2 online & 0.25011 & 0.25022 & 2.00 \\
Null & Fixed 26 online & 0.25086 & 0.25083 & 26.00 \\
Null & Fixed 26 gated & 0.25168 & 0.25088 & 26.00 \\
Additive & Withheld residual & 0.13648 & 0.13700 & 3.67 \\
Additive & Released residual & 0.13191 & 0.13125 & 4.00 \\
Additive & Withheld random & 0.13579 & 0.13622 & 3.83 \\
Additive & Released random & 0.13185 & 0.13125 & 4.00 \\
Additive & Fixed 2 online & 0.14356 & 0.14358 & 2.00 \\
Additive & Fixed 26 online & 0.13149 & 0.13158 & 26.00 \\
Additive & Fixed 26 gated & 0.13264 & 0.13167 & 26.00 \\
Interaction & Withheld residual & 0.14803 & 0.14711 & 6.00 \\
Interaction & Released residual & 0.14430 & 0.14254 & 6.00 \\
Interaction & Withheld random & 0.14907 & 0.14833 & 5.92 \\
Interaction & Released random & 0.14441 & 0.14244 & 6.00 \\
Interaction & Fixed 2 online & 0.20965 & 0.20959 & 2.00 \\
Interaction & Fixed 26 online & 0.13755 & 0.13780 & 26.00 \\
Interaction & Fixed 26 gated & 0.14050 & 0.13899 & 26.00 \\
Radial & Withheld residual & 0.12257 & 0.11954 & 3.00 \\
Radial & Released residual & 0.12143 & 0.11829 & 3.00 \\
Radial & Withheld random & 0.12207 & 0.11951 & 3.00 \\
Radial & Released random & 0.12087 & 0.11828 & 3.00 \\
Radial & Fixed 2 online & 0.19217 & 0.19210 & 2.00 \\
Radial & Fixed 26 online & 0.11714 & 0.11685 & 26.00 \\
Radial & Fixed 26 gated & 0.11867 & 0.11731 & 26.00 \\
\bottomrule
\end{tabular}

\caption{Complete mean outcomes of all seven arms. Tail risk follows a separate
ungated training budget and is outside the checkpoint guarantee. Values are
generated directly from saved results; uncertainty and paired comparisons are
reported above and in the raw CSV summaries.}
\end{table}
